% ================= ABSTRACT (INDONESIAN) =================
\chapter*{\large Judul Tugas Akhir, Format Times New Roman, Ukuran huruf 12}
\begin{center}
    \textbf{Abstrak}
\end{center}
\begin{singlespace}
Abstrak merupakan intisari dari skripsi dan ditulis dalam bahasa Indonesia. Abstrak disusun dalam satu paragraf dan panjangnya tidak lebih dari 250 kata yang diketik satu spasi. Abstrak hanya memuat isi penelitian (tujuan, metode, hasil, dan kesimpulan), tidak ada pengacuan pada pustaka, gambar dan tabel. Abstrak diketik dengan spasi satu, termasuk judul. Abstrak juga memuat kata kunci diletakkan di bagian paling bawah. Kata “abstrak” ditulis dengan huruf awal kapital dan diletakkan di tengah. Abstrak berbahasa Inggris (Abstract) diletakkan setelah halaman abstrak berbahasa Indonesia. Isi abstrak, ditulis dengan huruf Times New Roman ukuran 12, spasi 1. Bagian awal paragraph ditulis menjorok ke dalam 1 cm, dengan tulisan keseluruhan rata pinggir.

\vspace{0.5cm}
\noindent \textbf{Kata kunci:} ditulis dengan jenis dan ukuran huruf yang sama dengan abstrak.
\end{singlespace}
